% Generated draft: requires mathematical and source review.
\documentclass[11pt]{article}
\usepackage[margin=25mm]{geometry}
\usepackage{fontspec}
\setmainfont{Latin Modern Roman}
\usepackage{amsmath,amssymb,mathtools,unicode-math}
\setmathfont{Latin Modern Math}
\usepackage{enumitem,needspace,xurl,hyperref}
\setlength{\parindent}{0pt}
\setlength{\parskip}{5pt}
\setlength{\emergencystretch}{3em}
\newcommand{\N}{\mathbb{N}}
\newcommand{\Z}{\mathbb{Z}}
\newcommand{\Q}{\mathbb{Q}}
\newcommand{\R}{\mathbb{R}}
\newcommand{\C}{\mathbb{C}}
\newcommand{\F}{\mathbb{F}}
\newcommand{\A}{\mathbb{A}}
\newcommand{\PP}{\mathbb{P}}
\newcommand{\E}{\mathbb{E}}
\newcommand{\eps}{\varepsilon}
\newcommand{\dd}{\,\mathrm{d}}
\begin{document}
\section*{UnsolvedMath ID: 30003490}
\textbf{Source problem code:} catalogue entry
\subsection{Definitions and mathematical statement}
A finite discrete interpretation can be formulated as follows; the supplied record does not specify that this is its complete intended scope. Fix a positive integer \(d\), positive integers \(n_1,\ldots,n_d\), and a set of observed positions
\[
E\subseteq I:=\prod_{k=1}^{d}\{1,\ldots,n_k\}.
\]
A nonnegative partial tensor with observation pattern \(E\) is a vector \(x=(x_i)_{i\in E}\in\R_{\geq0}^{E}\). An independent probability completion of \(x\) consists of vectors \(p^{(k)}\in\R_{\geq0}^{n_k}\) satisfying
\[
\sum_{j=1}^{n_k}p^{(k)}_j=1\quad(1\leq k\leq d),
\qquad
x_{(i_1,\ldots,i_d)}=\prod_{k=1}^{d}p^{(k)}_{i_k}
\quad\bigl((i_1,\ldots,i_d)\in E\bigr).
\]
Thus the full tensor \(T_{i_1,\ldots,i_d}=\prod_{k=1}^{d}p^{(k)}_{i_k}\) is the joint probability mass function of independent random variables. Define
\[
\mathcal C_E=\left\{x\in\R_{\geq0}^{E}:x\text{ has an independent probability completion}\right\},
\qquad
\mathcal N_E=\R_{\geq0}^{E}\setminus\mathcal C_E.
\]
Under this interpretation, the question is whether, for every choice of \(d,n_1,\ldots,n_d,E\), the set \(\mathcal N_E\) is convex; explicitly, whether
\[
\forall x,y\in\mathcal N_E\ \forall\lambda\in[0,1],
\qquad \lambda x+(1-\lambda)y\in\mathcal N_E.
\]
This is a conditional formalization requiring review. The record does not specify the allowed tensor orders, dimensions, observation patterns, or ambient restrictions on the observed entries. For example, if admissible partial tensors must satisfy \(\sum_{i\in E}x_i\leq1\), the set whose convexity is asked about would instead be
\[
\left\{x\in\R_{\geq0}^{E}:\sum_{i\in E}x_i\leq1\right\}\setminus\mathcal C_E.
\]
The supplied evidence does not resolve which ambient set is intended. Convexity also requires a fixed observation pattern, or a specified common ambient space for comparing different patterns.
\subsection{Short English statement}
The question asks whether convex combinations of nonnegative partial tensors that admit no completion to an independent joint probability distribution must also admit no such completion. The displayed formulation uses finite discrete variables and a fixed observation pattern; the intended admissible entries and range of patterns remain unspecified.
\subsection{Sources}
\begin{itemize}
\item Supplied local registry record: permanent UnsolvedMath ID \(30003490\), title ``Convexity of Noncompletable Nonnegative Partial Tensors,'' exact source problem code \texttt{catalogue entry}. The record gives proposed year 2017 and names no proposer.
\item Supplied source citation: \emph{Algebraic Statistics (2018)}, Oberwolfach Reports, pp.\ 1207--1279, DOI \texttt{10.4171/owr/2017/20}. \url{https://doi.org/10.4171/owr/2017/20}. The supplied literature-source label is ``Oberwolfach Reports 2018.'' The background describes these reports as documenting mathematical workshops and open-problem discussions.
\item Supplied literature reference: ``Kahle et al., rank-one tensor completion,'' DOI \texttt{10.1137/16M1074102}. \url{https://doi.org/10.1137/16M1074102}. No further bibliographic details or specific result location are supplied.
\item Literature assessment checked at \texttt{2026-08-21}: the supplied curation labels the problem ``partially solved'' and reports that the proposed convexity holds in computed examples and in the diagonal case, while no proof or counterexample is known for the general partial-tensor set described in the record. This is supplied dated curation derived from the OWR record, not an independently verified assessment; its status and cited claims require independent verification.
\item The supplied statement status is \texttt{corrected\_verified}. Its verification note states that the official Oberwolfach report was checked directly and that the clean statement preserves the target while restoring definitions and removing source-dependent wording. This is a supplied verification claim; no independent source inspection is asserted here.
\end{itemize}
\end{document}
